\documentclass[11pt]{article}
\usepackage[utf8]{inputenc}
\usepackage[english]{babel}
\usepackage{graphicx}
\usepackage{hyperref}

\title{Carolus: Vision-Based Relative Navigation for a Rooted RoboMaster S1}
\author{} % TODO: author name(s)
\date{\today}

\begin{document}

\maketitle

\begin{abstract}
% TODO: content pending -- one-paragraph summary of the project, the
% approach (Carolus/UVGS-2 beacon + P4P, fused with wheel odometry via an
% EKF), and the main result obtained so far.
\end{abstract}

\section{Introduction and Setup}
% TODO: content pending -- hardware (rooted RoboMaster S1 via s1_sdk_hack,
% EP SDK, Raspberry Pi) and software setup (Carolus/UVGS-2: 4-LED beacon
% detection + P4P pose solve via Ceres).

\section{Methodology}
% TODO: content pending, by axis:
% - F3: robot_localization EKF (loosely-coupled, odom + Carolus absolute pose)
% - F4: MINS sandbox plan (KAIST dataset validation before any RoboMaster port)
% - F5: Carolus vs AprilTag orientation accuracy (hypothesis, not yet tested)
% - Ground-truth drift protocol (measuring-tape + CAD control points)

\section{Results to Date}
% TODO: content pending. IMPORTANT: the first F3 hardware test (2026-07-23,
% 1.6 m / 37 s) must be presented explicitly as a pre-test / sanity check,
% NOT as a drift-reduction demonstration -- see research-log/04-roadmap.md
% and research-log/07-perplexity/16-metrique-reduction-derive-hector.md.

\section{Planned Next Steps}
% TODO: content pending -- ground-truth drift protocol execution, MINS
% validation on KAIST, F5 test.

\section{Selected Lessons Learned}
% TODO: content pending -- structurally significant issues only, not the
% full internal bug registry.

\section{References}
% TODO: content pending -- prior intern documentation, supervisor's papers
% already reviewed (research-log/10-analyse-papiers-hector-mins-ekf.md).

\end{document}
